\section{Week 6 Refresher: Systems Arithmetic and Sparse Routing}

\subsection*{Setting the Scene}
The question shifts from ``can we train?'' to ``what is the bottleneck?'' Most of the math is accounting: FLOPs, bytes, bandwidth, and routing constraints.

\subsection*{Core Reminders}
\begin{itemize}[leftmargin=1.5em]
  \item \textbf{Throughput triangle:} compute limit, memory limit, communication limit.
  \item \textbf{MoE routing:} top-$k$ gate scores renormalized over selected experts.
  \item \textbf{Capacity limits:} overflow tokens can be dropped or rerouted.
  \item \textbf{Arithmetic intensity:} FLOPs per byte predicts whether kernels are compute- or memory-bound.
\end{itemize}

\subsection*{Pointers}
\begin{itemize}[leftmargin=1.5em]
  \item Optimization and memory accounting: see Week 2 refresher.
  \item Attention complexity and KV behavior: see Week 4 refresher.
  \item Decoding and token-rate implications: see Week 5 refresher.
\end{itemize}

\subsection*{Further Refresh}
Harchol-Balter queueing/performance texts, distributed systems lecture notes, MoE primary papers (Shazeer/Switch/Mixtral line).
